\documentclass[11pt]{article}
\newcommand{\videoid}{q9YDDtvI7fI}
% Preámbulo compartido de las notas del curso de topología algebraica.
% Cada nota empieza con \documentclass[11pt]{article}, define \videoid{...}
% (el identificador del video de YouTube) y luego hace % Preámbulo compartido de las notas del curso de topología algebraica.
% Cada nota empieza con \documentclass[11pt]{article}, define \videoid{...}
% (el identificador del video de YouTube) y luego hace % Preámbulo compartido de las notas del curso de topología algebraica.
% Cada nota empieza con \documentclass[11pt]{article}, define \videoid{...}
% (el identificador del video de YouTube) y luego hace \input{preambulo}.
% Compilar con LuaLaTeX:  lualatex nombre-de-la-nota.tex
\usepackage[spanish,provide=*]{babel}
\usepackage[T1]{fontenc}
\usepackage{lmodern}
\usepackage[a4paper,margin=2.6cm]{geometry}
\usepackage{amsmath,amssymb,amsthm}
\usepackage{xcolor}
\usepackage{tikz}
\usetikzlibrary{arrows.meta,babel}
\usepackage[colorlinks=true,urlcolor=blue!50!black,linkcolor=blue!50!black]{hyperref}

\hyphenpenalty=10000 \exhyphenpenalty=10000 \emergencystretch=3em

% Colores de las figuras (los mismos que usa la página)
\definecolor{extremo}{RGB}{47,107,79}
\definecolor{azul}{RGB}{44,74,154}

% Definiciones numeradas por unidad
\theoremstyle{definition}
\newtheorem{defn}{Definición}
% Numeración por unidad: \unidad{u}{n} (después de \input{preambulo}) hace que las
% definiciones de la nota se llamen u.n, u.(n+1), ...; n es el primer número
% libre de esa unidad (ver la tabla de CONVENCIONES.md).
\newcommand{\unidad}[2]{\renewcommand{\thedefn}{#1.\arabic{defn}}\setcounter{defn}{\numexpr#2-1\relax}}

% Observaciones y erratas del video. Usar SOLO si hay algo que señalar:
%   \begin{erratas}
%     \errata{203}{3:23}{Se dice ..., pero debe decir ...}
%   \end{erratas}
% El primer argumento son los segundos (enlaza al minuto exacto del video),
% el segundo es la marca visible, el tercero es el texto.
\newenvironment{erratas}{\section*{Observaciones y erratas del video}\begin{itemize}}{\end{itemize}}
\newcommand{\errata}[3]{%
  \edef\errurl{https://www.youtube.com/watch?v=\videoid\string&t=#1}%
  \item[\textbf{\expandafter\href\expandafter{\errurl}{#2}}] #3}
.
% Compilar con LuaLaTeX:  lualatex nombre-de-la-nota.tex
\usepackage[spanish,provide=*]{babel}
\usepackage[T1]{fontenc}
\usepackage{lmodern}
\usepackage[a4paper,margin=2.6cm]{geometry}
\usepackage{amsmath,amssymb,amsthm}
\usepackage{xcolor}
\usepackage{tikz}
\usetikzlibrary{arrows.meta,babel}
\usepackage[colorlinks=true,urlcolor=blue!50!black,linkcolor=blue!50!black]{hyperref}

\hyphenpenalty=10000 \exhyphenpenalty=10000 \emergencystretch=3em

% Colores de las figuras (los mismos que usa la página)
\definecolor{extremo}{RGB}{47,107,79}
\definecolor{azul}{RGB}{44,74,154}

% Definiciones numeradas por unidad
\theoremstyle{definition}
\newtheorem{defn}{Definición}
% Numeración por unidad: \unidad{u}{n} (después de % Preámbulo compartido de las notas del curso de topología algebraica.
% Cada nota empieza con \documentclass[11pt]{article}, define \videoid{...}
% (el identificador del video de YouTube) y luego hace \input{preambulo}.
% Compilar con LuaLaTeX:  lualatex nombre-de-la-nota.tex
\usepackage[spanish,provide=*]{babel}
\usepackage[T1]{fontenc}
\usepackage{lmodern}
\usepackage[a4paper,margin=2.6cm]{geometry}
\usepackage{amsmath,amssymb,amsthm}
\usepackage{xcolor}
\usepackage{tikz}
\usetikzlibrary{arrows.meta,babel}
\usepackage[colorlinks=true,urlcolor=blue!50!black,linkcolor=blue!50!black]{hyperref}

\hyphenpenalty=10000 \exhyphenpenalty=10000 \emergencystretch=3em

% Colores de las figuras (los mismos que usa la página)
\definecolor{extremo}{RGB}{47,107,79}
\definecolor{azul}{RGB}{44,74,154}

% Definiciones numeradas por unidad
\theoremstyle{definition}
\newtheorem{defn}{Definición}
% Numeración por unidad: \unidad{u}{n} (después de \input{preambulo}) hace que las
% definiciones de la nota se llamen u.n, u.(n+1), ...; n es el primer número
% libre de esa unidad (ver la tabla de CONVENCIONES.md).
\newcommand{\unidad}[2]{\renewcommand{\thedefn}{#1.\arabic{defn}}\setcounter{defn}{\numexpr#2-1\relax}}

% Observaciones y erratas del video. Usar SOLO si hay algo que señalar:
%   \begin{erratas}
%     \errata{203}{3:23}{Se dice ..., pero debe decir ...}
%   \end{erratas}
% El primer argumento son los segundos (enlaza al minuto exacto del video),
% el segundo es la marca visible, el tercero es el texto.
\newenvironment{erratas}{\section*{Observaciones y erratas del video}\begin{itemize}}{\end{itemize}}
\newcommand{\errata}[3]{%
  \edef\errurl{https://www.youtube.com/watch?v=\videoid\string&t=#1}%
  \item[\textbf{\expandafter\href\expandafter{\errurl}{#2}}] #3}
) hace que las
% definiciones de la nota se llamen u.n, u.(n+1), ...; n es el primer número
% libre de esa unidad (ver la tabla de CONVENCIONES.md).
\newcommand{\unidad}[2]{\renewcommand{\thedefn}{#1.\arabic{defn}}\setcounter{defn}{\numexpr#2-1\relax}}

% Observaciones y erratas del video. Usar SOLO si hay algo que señalar:
%   \begin{erratas}
%     \errata{203}{3:23}{Se dice ..., pero debe decir ...}
%   \end{erratas}
% El primer argumento son los segundos (enlaza al minuto exacto del video),
% el segundo es la marca visible, el tercero es el texto.
\newenvironment{erratas}{\section*{Observaciones y erratas del video}\begin{itemize}}{\end{itemize}}
\newcommand{\errata}[3]{%
  \edef\errurl{https://www.youtube.com/watch?v=\videoid\string&t=#1}%
  \item[\textbf{\expandafter\href\expandafter{\errurl}{#2}}] #3}
.
% Compilar con LuaLaTeX:  lualatex nombre-de-la-nota.tex
\usepackage[spanish,provide=*]{babel}
\usepackage[T1]{fontenc}
\usepackage{lmodern}
\usepackage[a4paper,margin=2.6cm]{geometry}
\usepackage{amsmath,amssymb,amsthm}
\usepackage{xcolor}
\usepackage{tikz}
\usetikzlibrary{arrows.meta,babel}
\usepackage[colorlinks=true,urlcolor=blue!50!black,linkcolor=blue!50!black]{hyperref}

\hyphenpenalty=10000 \exhyphenpenalty=10000 \emergencystretch=3em

% Colores de las figuras (los mismos que usa la página)
\definecolor{extremo}{RGB}{47,107,79}
\definecolor{azul}{RGB}{44,74,154}

% Definiciones numeradas por unidad
\theoremstyle{definition}
\newtheorem{defn}{Definición}
% Numeración por unidad: \unidad{u}{n} (después de % Preámbulo compartido de las notas del curso de topología algebraica.
% Cada nota empieza con \documentclass[11pt]{article}, define \videoid{...}
% (el identificador del video de YouTube) y luego hace % Preámbulo compartido de las notas del curso de topología algebraica.
% Cada nota empieza con \documentclass[11pt]{article}, define \videoid{...}
% (el identificador del video de YouTube) y luego hace \input{preambulo}.
% Compilar con LuaLaTeX:  lualatex nombre-de-la-nota.tex
\usepackage[spanish,provide=*]{babel}
\usepackage[T1]{fontenc}
\usepackage{lmodern}
\usepackage[a4paper,margin=2.6cm]{geometry}
\usepackage{amsmath,amssymb,amsthm}
\usepackage{xcolor}
\usepackage{tikz}
\usetikzlibrary{arrows.meta,babel}
\usepackage[colorlinks=true,urlcolor=blue!50!black,linkcolor=blue!50!black]{hyperref}

\hyphenpenalty=10000 \exhyphenpenalty=10000 \emergencystretch=3em

% Colores de las figuras (los mismos que usa la página)
\definecolor{extremo}{RGB}{47,107,79}
\definecolor{azul}{RGB}{44,74,154}

% Definiciones numeradas por unidad
\theoremstyle{definition}
\newtheorem{defn}{Definición}
% Numeración por unidad: \unidad{u}{n} (después de \input{preambulo}) hace que las
% definiciones de la nota se llamen u.n, u.(n+1), ...; n es el primer número
% libre de esa unidad (ver la tabla de CONVENCIONES.md).
\newcommand{\unidad}[2]{\renewcommand{\thedefn}{#1.\arabic{defn}}\setcounter{defn}{\numexpr#2-1\relax}}

% Observaciones y erratas del video. Usar SOLO si hay algo que señalar:
%   \begin{erratas}
%     \errata{203}{3:23}{Se dice ..., pero debe decir ...}
%   \end{erratas}
% El primer argumento son los segundos (enlaza al minuto exacto del video),
% el segundo es la marca visible, el tercero es el texto.
\newenvironment{erratas}{\section*{Observaciones y erratas del video}\begin{itemize}}{\end{itemize}}
\newcommand{\errata}[3]{%
  \edef\errurl{https://www.youtube.com/watch?v=\videoid\string&t=#1}%
  \item[\textbf{\expandafter\href\expandafter{\errurl}{#2}}] #3}
.
% Compilar con LuaLaTeX:  lualatex nombre-de-la-nota.tex
\usepackage[spanish,provide=*]{babel}
\usepackage[T1]{fontenc}
\usepackage{lmodern}
\usepackage[a4paper,margin=2.6cm]{geometry}
\usepackage{amsmath,amssymb,amsthm}
\usepackage{xcolor}
\usepackage{tikz}
\usetikzlibrary{arrows.meta,babel}
\usepackage[colorlinks=true,urlcolor=blue!50!black,linkcolor=blue!50!black]{hyperref}

\hyphenpenalty=10000 \exhyphenpenalty=10000 \emergencystretch=3em

% Colores de las figuras (los mismos que usa la página)
\definecolor{extremo}{RGB}{47,107,79}
\definecolor{azul}{RGB}{44,74,154}

% Definiciones numeradas por unidad
\theoremstyle{definition}
\newtheorem{defn}{Definición}
% Numeración por unidad: \unidad{u}{n} (después de % Preámbulo compartido de las notas del curso de topología algebraica.
% Cada nota empieza con \documentclass[11pt]{article}, define \videoid{...}
% (el identificador del video de YouTube) y luego hace \input{preambulo}.
% Compilar con LuaLaTeX:  lualatex nombre-de-la-nota.tex
\usepackage[spanish,provide=*]{babel}
\usepackage[T1]{fontenc}
\usepackage{lmodern}
\usepackage[a4paper,margin=2.6cm]{geometry}
\usepackage{amsmath,amssymb,amsthm}
\usepackage{xcolor}
\usepackage{tikz}
\usetikzlibrary{arrows.meta,babel}
\usepackage[colorlinks=true,urlcolor=blue!50!black,linkcolor=blue!50!black]{hyperref}

\hyphenpenalty=10000 \exhyphenpenalty=10000 \emergencystretch=3em

% Colores de las figuras (los mismos que usa la página)
\definecolor{extremo}{RGB}{47,107,79}
\definecolor{azul}{RGB}{44,74,154}

% Definiciones numeradas por unidad
\theoremstyle{definition}
\newtheorem{defn}{Definición}
% Numeración por unidad: \unidad{u}{n} (después de \input{preambulo}) hace que las
% definiciones de la nota se llamen u.n, u.(n+1), ...; n es el primer número
% libre de esa unidad (ver la tabla de CONVENCIONES.md).
\newcommand{\unidad}[2]{\renewcommand{\thedefn}{#1.\arabic{defn}}\setcounter{defn}{\numexpr#2-1\relax}}

% Observaciones y erratas del video. Usar SOLO si hay algo que señalar:
%   \begin{erratas}
%     \errata{203}{3:23}{Se dice ..., pero debe decir ...}
%   \end{erratas}
% El primer argumento son los segundos (enlaza al minuto exacto del video),
% el segundo es la marca visible, el tercero es el texto.
\newenvironment{erratas}{\section*{Observaciones y erratas del video}\begin{itemize}}{\end{itemize}}
\newcommand{\errata}[3]{%
  \edef\errurl{https://www.youtube.com/watch?v=\videoid\string&t=#1}%
  \item[\textbf{\expandafter\href\expandafter{\errurl}{#2}}] #3}
) hace que las
% definiciones de la nota se llamen u.n, u.(n+1), ...; n es el primer número
% libre de esa unidad (ver la tabla de CONVENCIONES.md).
\newcommand{\unidad}[2]{\renewcommand{\thedefn}{#1.\arabic{defn}}\setcounter{defn}{\numexpr#2-1\relax}}

% Observaciones y erratas del video. Usar SOLO si hay algo que señalar:
%   \begin{erratas}
%     \errata{203}{3:23}{Se dice ..., pero debe decir ...}
%   \end{erratas}
% El primer argumento son los segundos (enlaza al minuto exacto del video),
% el segundo es la marca visible, el tercero es el texto.
\newenvironment{erratas}{\section*{Observaciones y erratas del video}\begin{itemize}}{\end{itemize}}
\newcommand{\errata}[3]{%
  \edef\errurl{https://www.youtube.com/watch?v=\videoid\string&t=#1}%
  \item[\textbf{\expandafter\href\expandafter{\errurl}{#2}}] #3}
) hace que las
% definiciones de la nota se llamen u.n, u.(n+1), ...; n es el primer número
% libre de esa unidad (ver la tabla de CONVENCIONES.md).
\newcommand{\unidad}[2]{\renewcommand{\thedefn}{#1.\arabic{defn}}\setcounter{defn}{\numexpr#2-1\relax}}

% Observaciones y erratas del video. Usar SOLO si hay algo que señalar:
%   \begin{erratas}
%     \errata{203}{3:23}{Se dice ..., pero debe decir ...}
%   \end{erratas}
% El primer argumento son los segundos (enlaza al minuto exacto del video),
% el segundo es la marca visible, el tercero es el texto.
\newenvironment{erratas}{\section*{Observaciones y erratas del video}\begin{itemize}}{\end{itemize}}
\newcommand{\errata}[3]{%
  \edef\errurl{https://www.youtube.com/watch?v=\videoid\string&t=#1}%
  \item[\textbf{\expandafter\href\expandafter{\errurl}{#2}}] #3}

\unidad{1}{11} % unidad 1; primer enunciado: 1.11

\title{El grupo fundamental es, en efecto, un grupo}
\author{Luis Jorge Sánchez Saldaña\\[2pt]\small Facultad de Ciencias, UNAM}
\date{}

\begin{document}
\maketitle

\noindent\small Notas del video
\href{https://www.youtube.com/watch?v=q9YDDtvI7fI}{\guillemotleft{}El grupo fundamental es, en efecto, un grupo\guillemotright{}}
(curso de Introducción a la Topología Algebraica). Definen el producto en el grupo fundamental (la concatenación de lazos) y demuestran que, con esa operación, es un grupo.
\normalsize

\section*{El producto en el grupo fundamental}

Este video continúa el anterior. Allí se definió el grupo fundamental solo como conjunto: es el conjunto de clases de equivalencia de lazos basados en $x_0$, módulo la homotopía relativa a los extremos,
\[
  \pi_1(X,x_0)=\Omega(X,x_0)\big/\!\simeq\ \mathrm{rel}\ \{0,1\}.
\]
Se le llama \guillemotleft{}grupo\guillemotright{}, pero todavía no se ha dado una operación (más o menos sí, porque por eso se habló primero de la concatenación de caminos). La operación será concatenar lazos, lo cual tiene sentido: dos lazos basados en $x_0$ siempre se pueden concatenar, pues uno termina donde comienza el otro.

\medskip\noindent\textbf{Observación.} Si $[f],[g]\in\pi_1(X,x_0)$, se pueden tomar representantes $f$ y $g$, que son lazos basados en $x_0$, y concatenarlos. Como $f$ termina en $x_0$ y $g$ comienza en $x_0$, la concatenación $f\cdot g$ tiene sentido; además comienza y termina en $x_0$, así que es de nuevo un lazo, y su clase $[f\cdot g]$ es un elemento de $\pi_1(X,x_0)$. Es decir, la concatenación da un producto en el grupo fundamental.

\begin{teo}
$\pi_1(X,x_0)$ es un grupo con la operación
\[
  [f][g]=[f\cdot g].
\]
\end{teo}

El resto del video se dedica a demostrar que esta operación satisface los axiomas de grupo; por eso a $\pi_1(X,x_0)$ se le llama el grupo fundamental de $X$ basado en $x_0$.

\begin{proof}
Hay que ver que el producto está bien definido, que es asociativo, que existe un elemento neutro y que existen inversos.

\medskip\noindent\textit{El producto está bien definido.}
Como se multiplican clases de equivalencia usando representantes, hay que verificar que el resultado no depende de los representantes elegidos. Esto es consecuencia de la Proposición~1.8 del video anterior: si $[f]=[f_1]$ y $[g]=[g_1]$, entonces $[f\cdot g]=[f_1\cdot g_1]$. Además, como se vio arriba, el producto de dos lazos es un lazo.

\medskip\noindent\textit{Asociatividad.}
Se quiere ver que
\[
  \bigl([f][g]\bigr)[h]=[f]\bigl([g][h]\bigr),
\]
donde $f$, $g$ y $h$ son lazos arbitrarios basados en $x_0$. Por definición, esto es $[(f\cdot g)\cdot h]=[f\cdot(g\cdot h)]$.

En el dibujo los dos lazos se ven iguales, pero los lazos no son solo su imagen: son funciones parametrizadas, y para que dos lazos sean iguales tienen que coincidir como funciones. Hay que tener cuidado con lo que significan las concatenaciones. En $(f\cdot g)\cdot h$, en la primera mitad del tiempo se recorre $f\cdot g$ y en la segunda mitad se recorre $h$; como $f\cdot g$ se recorre en media unidad de tiempo, $f$ se recorre en un cuarto de tiempo y $g$ en otro cuarto. En $f\cdot(g\cdot h)$, en la primera mitad del tiempo se recorre $f$ y en la segunda mitad $g\cdot h$, así que $g$ y $h$ se recorren cada uno en un cuarto de tiempo (figura). Las imágenes y los recorridos coinciden, pero las parametrizaciones son distintas (por ejemplo, en el primero $f$ se recorre muy rápido y en el segundo no); por lo tanto, los lazos no son iguales.

\begin{center}
\begin{tikzpicture}[line cap=round]
  \def\W{8}
  % (f.g).h
  \node[left] at (-0.2,1.3) {$(f\cdot g)\cdot h$};
  \draw[very thick] (0,1.3) -- ({\W/4},1.3);
  \draw[very thick,azul] ({\W/4},1.3) -- ({\W/2},1.3);
  \draw[very thick,gray] ({\W/2},1.3) -- (\W,1.3);
  \node at ({\W/8},1.65) {\small $f$}; \node at ({3*\W/8},1.65) {\small $g$}; \node at ({3*\W/4},1.65) {\small $h$};
  % f.(g.h)
  \node[left] at (-0.2,0) {$f\cdot(g\cdot h)$};
  \draw[very thick] (0,0) -- ({\W/2},0);
  \draw[very thick,azul] ({\W/2},0) -- ({3*\W/4},0);
  \draw[very thick,gray] ({3*\W/4},0) -- (\W,0);
  \node at ({\W/4},0.35) {\small $f$}; \node at ({5*\W/8},0.35) {\small $g$}; \node at ({7*\W/8},0.35) {\small $h$};
  % marcas del intervalo
  \foreach \x/\l in {0/0,2/{1/4},4/{1/2},6/{3/4},8/1}{\draw (\x,-0.15) -- (\x,-0.05); \node[below,gray] at (\x,-0.15) {\small $\l$};}
\end{tikzpicture}
\par\smallskip
\begin{minipage}{0.9\textwidth}\small\centering
Las dos concatenaciones recorren los mismos tres lazos, pero en tiempos distintos: la primera dedica un cuarto del tiempo a $f$, otro a $g$ y la mitad a $h$; la segunda, la mitad a $f$ y un cuarto a cada uno de $g$ y $h$.
\end{minipage}
\end{center}

Lo que sí sucede es que estos dos lazos son homotópicos relativos a los extremos, y con eso se obtiene la igualdad de las clases. La homotopía es fácil de imaginar: se reparametriza, estirando el pedazo de $f$ hasta que ocupe la mitad del intervalo, desplazando el pedazo de $g$ y achicando el de $h$ hasta que quede de tamaño un cuarto. Formalmente, se quiere ver que $(f\cdot g)\cdot h\simeq f\cdot(g\cdot h)\ \mathrm{rel}\ \{0,1\}$; para ello se define $F\colon I\times I\to X$ como
\[
  F(s,t)=
  \begin{cases}
    f\!\left(\dfrac{4s}{t+1}\right) & \text{si } s\in\bigl[0,\tfrac{t+1}{4}\bigr],\\[8pt]
    g(4s-t-1) & \text{si } s\in\bigl[\tfrac{t+1}{4},\tfrac{t+2}{4}\bigr],\\[8pt]
    h\!\left(\dfrac{4s-t-2}{2-t}\right) & \text{si } s\in\bigl[\tfrac{t+2}{4},1\bigr].
  \end{cases}
\]
\begin{center}
\begin{tikzpicture}[line cap=round]
  \draw[very thick] (0,3.6) -- (3.6,3.6);
  \draw[very thick,azul] (0,0.0) -- (3.6,0.0);
  \draw[very thick,extremo] (0,0.0) -- (0,3.6);
  \draw[very thick,extremo] (3.6,0.0) -- (3.6,3.6);
  \draw[thick,gray] (0.9,3.6) -- (1.8,0.0);
  \draw[thick,gray] (1.8,3.6) -- (2.7,0.0);
  \node[left] at (-0.1,3.6) {$(f\cdot g)\cdot h$};
  \node[left,azul] at (-0.1,0.0) {$f\cdot(g\cdot h)$};
  \node[above] at (0.45,3.6) {\small $f$};
  \node[above] at (1.35,3.6) {\small $g$};
  \node[above] at (2.7,3.6) {\small $h$};
  \node[below,azul] at (0.9,0.0) {\small $f$};
  \node[below,azul] at (2.25,0.0) {\small $g$};
  \node[below,azul] at (3.15,0.0) {\small $h$};
  \node[gray!80!black] at (0.504,1.8) {\small $f$};
  \node[gray!80!black] at (1.8,1.8) {\small $g$};
  \node[gray!80!black] at (3.024,1.8) {\small $h$};
\end{tikzpicture}
\par\smallskip
\begin{minipage}{0.9\textwidth}\small\centering
Cuadrado de la homotopía de asociatividad: arriba ($t=0$) está $(f\cdot g)\cdot h$ y abajo ($t=1$) está $f\cdot(g\cdot h)$. Las dos diagonales siguen cómo se desplazan los puntos donde terminan $f$ y $g$: de $\tfrac14$ y $\tfrac12$ arriba a $\tfrac12$ y $\tfrac34$ abajo. Los lados verticales (en verde) van a $x_0$.
\end{minipage}
\end{center}

El video deja como un ejercicio muy divertido verificar que $F$ es, en efecto, una homotopía relativa a los extremos entre esos dos lazos.

\medskip\noindent\textit{Elemento neutro.}
Se quiere ver que existe un elemento $[c]\in\pi_1(X,x_0)$ tal que
\[
  [f][c]=[c][f]=[f]\qquad\text{para todo } [f]\in\pi_1(X,x_0).
\]
Este lazo misterioso es el \emph{lazo constante} $c\colon I\to X$, dado por $c(s)=x_0$ para todo $s\in I$: comienza en $x_0$, nunca se mueve y termina en $x_0$. El video verifica solo una de las igualdades, $[f\cdot c]=[f]$, es decir, que $f\cdot c\simeq f\ \mathrm{rel}\ \{0,1\}$. La idea es similar a la anterior: se toma el pedazo de $f$ y se estira hasta que ocupe todo el intervalo, mientras el lazo constante queda completamente comprimido a un punto del lado derecho; siempre se respetan los extremos. Formalmente, la homotopía $F\colon I\times I\to X$ es
\[
  F(s,t)=
  \begin{cases}
    f\!\left(\dfrac{2s}{t+1}\right) & \text{si } s\in\bigl[0,\tfrac{t+1}{2}\bigr],\\[8pt]
    x_0 & \text{si } s\in\bigl[\tfrac{t+1}{2},1\bigr].
  \end{cases}
\]
\begin{center}
\begin{tikzpicture}[line cap=round]
  \draw[very thick] (0,3.6) -- (3.6,3.6);
  \draw[very thick,azul] (0,0.0) -- (3.6,0.0);
  \draw[very thick,extremo] (0,0.0) -- (0,3.6);
  \draw[very thick,extremo] (3.6,0.0) -- (3.6,3.6);
  \draw[thick,gray] (1.8,3.6) -- (3.6,0.0);
  \node[left] at (-0.1,3.6) {$f\cdot c$};
  \node[left,azul] at (-0.1,0.0) {$f$};
  \node[above] at (0.9,3.6) {\small $f$};
  \node[above] at (2.7,3.6) {\small $c$};
  \node[below,azul] at (1.8,0.0) {\small $f$};
  \node[gray!80!black] at (1.08,1.62) {\small $f$};
  \node[gray!80!black] at (2.88,2.7) {\small $x_0$};
\end{tikzpicture}
\par\smallskip
\begin{minipage}{0.9\textwidth}\small\centering
Cuadrado de la homotopía del neutro: arriba ($t=0$) está $f\cdot c$ y abajo ($t=1$) está $f$. La diagonal sigue el punto donde termina $f$: va de $\tfrac12$ arriba a $1$ abajo, de modo que el lazo constante queda comprimido a un punto en el lado derecho. Los lados verticales (en verde) van a $x_0$.
\end{minipage}
\end{center}

Con esto $[c]$ es un neutro para la multiplicación.

\medskip\noindent\textit{Inversos.}
Se quiere ver que para cada $[f]\in\pi_1(X,x_0)$ existe $[\bar f]\in\pi_1(X,x_0)$ tal que
\[
  [f][\bar f]=[\bar f][f]=[c].
\]
La idea: $\bar f$ es el mismo lazo $f$ parametrizado al revés, con la misma imagen pero recorrido en la dirección opuesta. Entonces $f\cdot\bar f$ es ir y regresar por el mismo camino, y eso se puede contraer al lazo constante, como si se aventara una cuerda y luego se fuera recogiendo: cada vez se regresa un poco antes, hasta que solo queda $x_0$. Es esencial que en ningún momento se muevan los extremos; si se movieran, ya no se estaría respetando la relación de equivalencia.

Formalmente, se define $\bar f\colon I\to X$ por
\[
  \bar f(s)=f(1-s),
\]
y se define $F\colon I\times I\to X$ por
\[
  F(s,t)=
  \begin{cases}
    f(2s) & \text{si } 2s\in[0,t],\\[2pt]
    f(t) & \text{si } 2s\in[t,2-t],\\[2pt]
    \bar f(2s-1) & \text{si } 2s\in[2-t,2].
  \end{cases}
\]
\begin{center}
\begin{tikzpicture}[line cap=round]
  \draw[very thick] (0,3.6) -- (3.6,3.6);
  \draw[very thick,azul] (0,0.0) -- (3.6,0.0);
  \draw[very thick,extremo] (0,0.0) -- (0,3.6);
  \draw[very thick,extremo] (3.6,0.0) -- (3.6,3.6);
  \draw[thick,gray] (0.0,3.6) -- (1.8,0.0);
  \draw[thick,gray] (3.6,3.6) -- (1.8,0.0);
  \node[left] at (-0.1,3.6) {$c$};
  \node[left,azul] at (-0.1,0.0) {$f\cdot\bar f$};
  \node[below,azul] at (0.9,0.0) {\small $f$};
  \node[below,azul] at (2.7,0.0) {\small $\bar f$};
  \node[gray!80!black] at (0.504,1.08) {\small $f$};
  \node[gray!80!black] at (1.8,2.52) {\small $f(t)$};
  \node[gray!80!black] at (3.096,1.08) {\small $\bar f$};
\end{tikzpicture}
\par\smallskip
\begin{minipage}{0.9\textwidth}\small\centering
Cuadrado de la homotopía de los inversos: arriba ($t=0$) está el lazo constante $c$ y abajo ($t=1$) está $f\cdot\bar f$. Las dos diagonales parten de las esquinas de arriba y se juntan en el punto medio de abajo: a la izquierda se recorre $f$ cada vez más lejos, a la derecha $\bar f$, y en el triángulo del centro el lazo se queda quieto en $f(t)$. Los lados verticales (en verde) van a $x_0$.
\end{minipage}
\end{center}

Con esta fórmula, $F(s,0)=c(s)$ y $F(s,1)=(f\cdot\bar f)(s)$: la tapa de arriba va al lazo constante, la de abajo a la concatenación de $f$ con $\bar f$ (figura), y los lados del cuadrado van a $x_0$. El video invita a comprobar que $F$ es una homotopía relativa a los extremos entre esos dos lazos. Por lo tanto $f\cdot\bar f\simeq c\ \mathrm{rel}\ \{0,1\}$ (como $\simeq$ es simétrica, da lo mismo en qué sentido se lea), es decir, $[f][\bar f]=[c]$, y así existen inversos.

\medskip
Se demostró que el producto está bien definido, que es asociativo, que hay neutro y que hay inversos; por lo tanto $\pi_1(X,x_0)$ es un grupo.
\end{proof}

El video cierra señalando que ya se tiene el grupo fundamental, con el que se trabajará en los siguientes videos.

\begin{erratas}
  \errata{1045}{17:25}{Se dice que la tapa de arriba va a $f\cdot\bar f$ y la de abajo al lazo constante, pero con la fórmula escrita se tiene $F(s,0)=c$ y $F(s,1)=f\cdot\bar f$: arriba va el lazo constante y abajo $f\cdot\bar f$.}
\end{erratas}

\end{document}
